\section{Related Work}
\label{related_work}
\textbf{Commit and Pull Request Understanding}
Understanding and analyzing PRs/commits is a fundamental task in software engineering research, aiming to support tasks such as code review, bug triage, and change impact analysis~\cite{icsik2025enhancing,chen2022untangling,li2022utango,dias2015untangling,muylaert2018untangling,herbold2022fine,liu2019automatic,xiao2024generative,bai2025clg,meng2024keytitle,meijer2024ecosystem,venigalla2024exploratory,pan2022automated}.
Liu et al.\cite{liu2019automatic} proposed a sequence-to-sequence model with a pointer-generator mechanism to automatically generate pull request (PR) descriptions by leveraging commit messages and code comments. 
They further applied reinforcement learning to directly optimize ROUGE scores and mitigate the mismatch between training objectives and evaluation metrics.
Bai et al.~\cite{bai2025clg} introduced the CLG framework, which generates PR checklists automatically from contribution guidelines by combining multi-label classification and text generation techniques. 
Their method significantly outperformed baselines on summarization tasks and demonstrated practical value in open-source community management.
% Meng et al.~\cite{meng2024keytitle} developed KeyTitle, a method that formulates the generation of bug report titles as a single-sentence summarization task. 
% By integrating keyword planning with a Transformer-based encoder-decoder model, their approach achieved up to 30.4\% improvement in ROUGE scores on GitHub datasets.


\textbf{Security and Maintenance Challenges in Third-Party Library Dependencies.}
Prior work has extensively studied the risks posed by third-party libraries, including vulnerability propagation, delayed patch adoption, and ecosystem-wide dependency management challenges~\cite{zimmermann2019small,duan2020towards, huang2022characterizing,decan2019empirical,alfadel2023empirical,liu2022demystifying,chen2023identifying,zhang2023compatible,williams2025research,zhao2023software,zhang2023mitigating,zhang2024does,zhang2024vulnerability,gao2025vulnerability,zhou2024magneto,chen2025diffploitfacilitatingcrossversionexploit,liu2025empirical}. 
Zeng et al.~\cite{zeng2024survey} present a comprehensive survey on the security of third-party dependencies, highlighting trends in detection, mitigation, and secure integration practices.
Wu et al.~\cite{wu2023understanding} perform a large-scale empirical study over 44,450 triplets in the Maven ecosystem, revealing that 73.3\% of downstream projects depending on vulnerable libraries remain practically unaffected—underscoring the importance of precise impact assessment.
Chen et al.~\cite{chen2024exploiting} propose \textsc{Vesta} a vulnerability-oriented test generation framework that uses exploit code migration to evaluate real-world vulnerability reachability, improving detection success by 53.4\% over baselines.
Unlike prior work, \appname targets subtle, non-crashing defects in DL libraries by bridging defect modeling and downstream impact analysis through extracting fine-grained trigger conditions from code changes and semantically matching them to real client usage.
